\SourceLeaf{332}
$(\mathrm{mod.}\,770)$ his $(39,237)$, $(171,-27)$, $(269,-83)$, $(291,-127)$, $(-39,-237)$, $(-171,27)$, $(-269,83)$, $(-291,127)$, quatuor posteriores sunt reiiciendi, tamquam quatuor prioribus oppositi. Ceterum perspicuum est, si $(B,B')$ sit valor sibi ipsi oppositus, $2B$, $2B'$ et proin etiam $2\Delta p$, $2\Delta q$, $2\Delta r$ per $D$ divisibiles fore; quodsi itaque $\Delta$, $D$ inter se primi sunt, etiam $2p$, $2q$, $2r$ per $D$ divisibiles erunt. et quum, per I, in hoc casu etiam $p$, $q$, $r$ divisorem communem habere nequeant, etiam $2$ per $D$ divisibilis esse debebit, quod fieri nequit nisi $D$ vel $=\pm1$, vel $=\pm2$. Quamobrem pro omnibus valoribus ipsius $D$ maioribus quam $2$ semper erit $n=0$, si $\Delta$ ad $D$ est primus.

III. His ita factis manifestum est, quamvis repraesentationem propriam formae $\varphi$ per $f$ necessario ad aliquem e valoribus remanentibus pertinere debere, et quidem ad unicum tantum. Quare hi valores successive sunt percurrendi, repraesentationesque ad singulos pertinentes investigandae. Ut inveniantur repraesentationes ad valorem \textit{datum} $(B,B')$ pertinentes, primo determinanda est forma ternaria $g=\left(\begin{smallmatrix}a,\,a',\,a''\\b,\,b',\,b''\end{smallmatrix}\right)$, cuius determinans $=\Delta$ et in qua $a=p$, $b''=q$, $a'=r$, $ab-b'b''=B$, $a'b'-bb''=B'$; valores ipsorum $a''$, $b$, $b'$ hinc inveniuntur adiumento aequationum in II art.~276, ex quibus facile perspicitur, in eo casu, ubi $\Delta$, $D$ inter se primi sint, $b,b'\,a''$ necessario fieri integros (nempe quoniam hi tres numeri, multiplicati tum per $D$ tum per $\Delta$ integros producunt). Iam si vel aliquis coëfficientium $b,b',b''$ fractus est, vel formae $f$, $g$ non sunt aequivalentes: nullae repraesentationes formae $\varphi$ per $f$ ad $(B,B')$ pertinentes dari possunt; si vero $b,b',a''$ sunt integri, formaeque $f$, $g$ aequivalentes, quaevis transformatio illius in hanc, ut
\[
\begin{array}{ccc}
\alpha,&\beta,&\gamma\\
\alpha',&\beta',&\gamma'\\
\alpha'',&\beta'',&\gamma''
\end{array}
\]
talem repraesentationem suppeditat, puta
\[
x=\alpha t+\beta u,\quad x'=\alpha't+\beta'u,\quad x''=\alpha''t+\beta''u
\]
manifestoque nulla huiusmodi repraesentatio exstare poterit, quae non ex aliqua transformatione deduci posset. Hoc itaque modo ea problematis secundi pars, quae investigat repraesentationes \textit{proprias}, ad problema tertium iam est reducta.
\SourceLeaf{333}

IV. Ceterum transformationes diversae formae $f$ in $g$ semper producunt repraesentationes diversas, eo solo casu excepto, ubi valor $(B,B')$ sibi ipsi oppositus est, in quo binae transformationes unicam semper repraesentationem suppeditant. Supponende enim, $f$ transire in $g$ etiam per substitutionem
\[
\begin{array}{ccc}
\alpha,&\beta,&\delta\\
\alpha',&\beta',&\delta'\\
\alpha'',&\beta'',&\delta''
\end{array}
\]
(quae eandem repr. praebet ut transf. praec.), denotandoque per $k$, $\mathfrak{k}$, $\zeta$, $\eta$ numeros eosdem ut in II art. praec., erit
\[
B=k\mathfrak{k}B+\eta\mathfrak{k}D,\quad B'=k\mathfrak{k}B'+\zeta\mathfrak{k}D
\]
si itaque vel uterque $k$, $\mathfrak{k}$ supponitur $=+1$, vel uterque $=-1$, erit (quia casum $D=0$ exclusimus) $\zeta=0$, $\eta=0$, unde facile sequitur $\delta=\gamma$, $\delta'=\gamma'$, $\delta''=\gamma''$; quare illae duae transformationes in eo solo casu diversae esse possunt, ubi alter numerorum $k$, $\mathfrak{k}$ est $+1$, alter $-1$; tunc erit $B\equiv-B$, $B'\equiv-B'$, $(\mathrm{mod.}\,D)$, sive valor $(B,B')$ sibi ipsi oppositus.

V. Ex iis, quae supra (art.~271) de criteriis formarum definitarum et indefinitarum tradidimus, facile sequitur, si $\Delta$ sit positivus, $D$ negativus, atque $\varphi$ forma negativa, $g$ fieri formam definitam negativam; si vero $\Delta$ sit positivus, atque vel $D$ positivus, vel $D$ negativus et $\varphi$ forma positiva, $g$ evadere formam indefinitam. Iam quum $f$, $g$ certo aequivalentes esse nequeant, nisi respectu huius qualitatis similes sint, manifestum est, formas binarias determinantis positivi nec non positivas, per ternariam negativam proprie repraesentari non posse, neque formas binarias negativas per ternariam indefinitam determinantis positivi; sed per formam ternariam prioris posteriorisve speciei unice binarias posterioris priorisve resp. Simili modo concluditur, per formam ternariam determinantis negativi definitam (\textit{i.~e.} positivam) unice repraesentari binarias positivas, per indefinitam unice negativas et formas det. positivi.

\Needspace{8\baselineskip}
\Article{284}
Quum repraesentationes \textit{impropriae} formae binariae $\varphi$ determinantis $D$ per ternariam $f$, cui adiuncta est $F$, eae sint, ex quibus repraesentationes
\SourceLeaf{334}
impropriae numeri $D$ per formam $F$ sequuntur, $\varphi$ per $f$ manifesto nequit improprie repraesentari, nisi $D$ factores quadratos implicet. Ponamus, omnia quadrata ipsum $D$ metientia (praeter $1$) esse $ee$, $e'e'$, $e''e''$ etc. (quorum multitudo finita erit, quia supponimus, non esse $D=0$), praebebitque quaelibet repr. impr. formae $\varphi$ per $f$ repraesentationem numeri $D$ per $F$, in qua valores indeterminatarum aliquem e numeris $e$, $e'$, $e''$ etc. pro divisore communi maximo habent; hoc respectu brevitatis caussa quamvis repr. impr. formae $\varphi$ ad divisorem quadratum $ee$ vel $e'e'$ vel $e''e''$ etc. \textit{pertinere} dicemus. Iam omnes repr. formae $\varphi$ ad eundem divisorem quadratum \textit{datum} $ee$ (cuius radicem $e$ positive acceptam supponimus) pertinentes per regulas sequentes inveniuntur, ex quarum demonstratione synthetica, propter brevitatem hic praeferenda, analysis per quam evolutae sunt, facile restitui poterit.

\textit{Primo} eruantur omnes formae binariae determinantis $\frac{D}{ee}$, quae in formam $\varphi$ transeunt per substitutionem propriam talem $T=\kappa t+\lambda u$, $U=\mu u$, designantibus $T$, $U$ indeterminatas talis formae; $t$, $u$ indet. formae $\varphi$; $\kappa$, $\mu$ integros positivos (quorum productum itaque $=e$); $\lambda$ integrum positivum minorem quam $\mu$ (sive etiam cifram). Hae formae, cum transformationibus respondentibus, ita inveniuntur:

Aequetur $\kappa$ successive singulis divisoribus ipsius $e$ positive acceptis (inclusis etiam $1$ et $e$), fiatque $\mu=\frac{e}{\kappa}$; pro singulis valoribus determinatis ipsorum $\kappa$, $\mu$ tribuantur ipsi $\lambda$ omnes valores integri a $0$ usque ad $\mu-1$, quo pacto omnes transformationes certo habebuntur. Iam forma, quae per quamvis substitutionem $T=\kappa t+\lambda u$, $U=\mu u$ in $\varphi$ transit, invenitur investigando formam, in quam $\varphi$ transit per hanc $t=\frac{1}{\kappa}T-\frac{\lambda}{e}U$, $u=\frac{1}{\mu}U$; sic formae singulis transformationibus respondentes obtinebuntur; sed ex omnibus his formis eae tantum retinendae sunt, in quibus omnes tres coëfficientes evadunt integri\HistoricalNote{*)}{Si de hoc problemate fusius agere hic liceret, solutionem admodum contrahere possemus. Id statim obvium est, pro $\kappa$ alios divisores ipsius $e$ accipere non esse necessarium, nisi quorum quadratum metiatur coëfficientem primum formae $\varphi$. Ceterum hoc problema, ex quo etiam solutiones simpliciores probl. artt.~213, 214 deduci possunt, alia occasione idonea resumere nobis reservamus.}.

\textit{Secundo} ponamus $\Phi$ esse aliquam ex hisce formis, quae in $\varphi$ transeat per subst. $T=\kappa t+\lambda u$, $U=\mu u$; investigentur omnes repraesentationes \textit{propriae}
\SourceLeaf{335}
formae $\Phi$ per $f$ (si quae dantur), exhibeanturque indefinite per
\[
x=\mathfrak{A}T+\mathfrak{B}U,\quad x'=\mathfrak{A}'T+\mathfrak{B}'U,\quad x''=\mathfrak{A}''T+\mathfrak{B}''U\quad\ldots\quad(\mathfrak{R})
\]
denique ex singulis $(\mathfrak{R})$ deducatur repraesentatio
\[
x=\alpha t+\beta u,\quad x'=\alpha't+\beta'u,\quad x''=\alpha''t+\beta''u\quad\ldots\quad(\rho)
\]
per aequationes
\[
\begin{array}{llll}
\alpha=\kappa\mathfrak{A},&\alpha'=\kappa\mathfrak{A}',&\alpha''=\kappa\mathfrak{A}''&\ldots\quad(R)\\
\beta=\lambda\mathfrak{A}+\mu\mathfrak{B},&\beta'=\lambda\mathfrak{A}'+\mu\mathfrak{B}',&\beta''=\lambda\mathfrak{A}''+\mu\mathfrak{B}''&
\end{array}
\]
Eodem prorsus modo, ut forma $\Phi$, tractentur formae reliquae per regulam primam inventae (si plures adsunt), ita ut ex singulis cuiusque repraesentationibus propriis aliae repraesentationes deriventur, dicoque, hoc modo prodire cunctas repraesentationes formae $\varphi$ ad divisorem $ee$ pertinentes, et quidem quamlibet semel tantum.

\textit{Dem.} I. Formam ternariam $f$ per quamvis substitutionem $(\rho)$ revera transire in $\varphi$, tam obvium est, ut explicatione ampliori non opus sit; quamlibet autem repr. $(\rho)$ esse impropriam et ad divisorem $ee$ pertinere, inde patet, quod numeri $\alpha'\beta''-\alpha''\beta'$, $\alpha''\beta-\alpha\beta''$, $\alpha\beta'-\alpha'\beta$ resp. fiunt $=e(\mathfrak{A}'\mathfrak{B}''-\mathfrak{A}''\mathfrak{B}')$, $e(\mathfrak{A}''\mathfrak{B}-\mathfrak{A}\mathfrak{B}'')$, $e(\mathfrak{A}\mathfrak{B}'-\mathfrak{A}'\mathfrak{B})$, unde illorum divisor comm. max. manifesto erit $e$ (quoniam $(\mathfrak{R})$ est repraesentatio propria).

II. Ostendemus, ex quavis repraesentatione data $(\rho)$ formae $\varphi$, inveniri posse repraesentationem propriam formae determinantis $\frac{D}{ee}$, inter formas per regulam primam inventas contentae, sive ex valoribus datis ipsorum $\alpha,\allowbreak \alpha',\allowbreak \alpha'',\allowbreak \beta,\allowbreak \beta',\allowbreak \beta''$ deduci posse valores integros ipsorum $\kappa$, $\lambda$, $\mu$, conditionibus praescriptis, atque valores ipsorum $\mathfrak{A}$, $\mathfrak{A}'$, $\mathfrak{A}''$, $\mathfrak{B}$, $\mathfrak{B}'$, $\mathfrak{B}''$, aequationibus $(R)$ satisfacientes, et quidem unico tantum modo. Primo statim patet ex tribus aequ. primis in $(R)$, pro $\kappa$ accipi debere divisorem communem maximum ipsorum $\alpha$, $\alpha'$, $\alpha''$ signo positivo (quum enim $\mathfrak{A}'\mathfrak{B}''-\mathfrak{A}''\mathfrak{B}'$, $\mathfrak{A}''\mathfrak{B}-\mathfrak{A}\mathfrak{B}''$, $\mathfrak{A}\mathfrak{B}'-\mathfrak{A}'\mathfrak{B}$ divisorem communem non habere debeant, etiam $\mathfrak{A}$, $\mathfrak{A}'$, $\mathfrak{A}''$ div. comm. habere nequeunt); hinc etiam $\mathfrak{A}$, $\mathfrak{A}'$, $\mathfrak{A}''$ determinati erunt, nec non $\mu=\frac{e}{\kappa}$ (quem necessario integrum fieri facile perspicitur). Ponamus, tres integros $\mathfrak{a}$, $\mathfrak{a}'$, $\mathfrak{a}''$ ita acceptos esse, ut fiat $\mathfrak{a}\mathfrak{A}+\mathfrak{a}'\mathfrak{A}'+\mathfrak{a}''\mathfrak{A}''=1$, scribamusque brevitatis caussa $k$ pro $\mathfrak{a}\mathfrak{B}+\mathfrak{a}'\mathfrak{B}'+\mathfrak{a}''\mathfrak{B}''$.
\SourceLeaf{336}
Tunc ex tribus ultimis aeq. $(R)$ sequitur, esse debere $\mathfrak{a}\beta+\mathfrak{a}'\beta'+\mathfrak{a}''\beta''=\lambda+\mu k$, unde statim patet, pro $\lambda$ unicum tantummodo valorem inter limites $0$ et $\mu-1$ situm dari. Quo facto quum etiam $\mathfrak{B}$, $\mathfrak{B}'$, $\mathfrak{B}''$ valores determinatos nanciscantur. nihil superest, nisi ut demonstremus, hos semper hinc integros evadere. Fiet autem
\[
\begin{aligned}
\mathfrak{B}=\frac{1}{\mu}(\beta-\lambda\mathfrak{A})
&=\frac{1}{\mu}(\beta(1-\mathfrak{a}\mathfrak{A})-\mathfrak{A}(\mathfrak{a}'\beta'+\mathfrak{a}''\beta''))+\mathfrak{A}k\\
&=\frac{1}{\mu}(\mathfrak{a}''(\mathfrak{A}''\beta-\mathfrak{A}\beta'')-\mathfrak{a}'(\mathfrak{A}\beta'-\mathfrak{A}'\beta))+\mathfrak{A}k\\
&=\frac{1}{e}(\mathfrak{a}''(\alpha''\beta-\alpha\beta'')-\mathfrak{a}'(\alpha\beta'-\alpha'\beta))+\mathfrak{A}k
\end{aligned}
\]
eritque adeo manifesto integer, similiterque facile confirmatur, etiam ipsos $\mathfrak{B}'$, $\mathfrak{B}''$ valores integros nancisci. --- Ex his ratiociniis colligitur, nullam repraesentationem impropriam formae $\varphi$ per $f$, ad divisorem $ee$ pertinentem, exstare posse, quae per methodum traditam vel non vel pluries obtineatur.

Quodsi iam eodem modo reliqui divisores quadrati ipsius $D$ tractantur, repraesentationesque ad singulos pertinentes eruuntur, cunctae repraesentationes impropriae formae $\varphi$ per $f$ habebuntur.

Ceterum ex hac solutione facile deducitur, theorema ad finem art. praec. pro repraess. propriis traditum etiam ad improprias patere, scilicet generaliter nullam formam binariam positivam det. negativi per ternariam negativam repraesentari posse etc.; patet enim, si $\varphi$ sit forma talis binaria, quae propter illud theorema per $f$ proprie repraesentari nequeat, etiam omnes formas determinantium $\frac{D}{ee}$, $\frac{D}{e'e'}$ etc., ipsam $\varphi$ implicantes per $f$ proprie repraesentari non posse, quum hae formae omnes determinantem eodem signo affectum habeant ut $\varphi$. et, quoties hi determinantes negativi sunt, vel omnes evadant formae positivae vel negativae, prout $\varphi$ ad illas vel ad has pertinet.

\Needspace{8\baselineskip}
\Article{285}
De quaestionibus problema tertium nobis propositum constituentibus (ad quod duo priora in praecc. sunt reducta), scilicet propositis duabus formis ternariis eiusdem determinantis, diiudicare, utrum aequivalentes sint necne, et in casu priori omnes transformationes alterius in alteram invenire, pauca tantum hoc loco inserere possumus, quum solutio completa, qualem pro problematibus analogis in formis binariis tradidimus, hic adhuc maioribus difficultatibus sit obnoxia.
\SourceLeaf{337}
Quamobrem ad quosdam casus particulares, propter quos praecipue haecce digressio instituta est, disquisitionem nostram limitabimus.

I\quad Pro determinante $+1$ supra ostensum est, omnes formas ternarias in duas classes distribui, quarum altera omnes formas indefinitas, altera omnes definitas (negativas) contineat. Hinc statim concluditur, duas formas ternarias quascunque det.~$1$ aequivalentes esse, si vel utraque sit definita vel utraque indefinita; si vero altera sit definita, altera indefinita, aequivalentiam locum non habere (propositionis pars posterior manifesto valet generaliter pro formis determinantis cuiuscunque). --- Simili modo duae formae quaecunque determinantis $-1$ certo aequivalebunt, si vel utraque definita est, vel utraque indefinita. --- Duae formae definitae determinantis $2$ semper aequivalebunt; duae indefinitae non aequivalebunt, si in altera tres coëfficientes primi omnes pares sunt, in altera vero non omnes sunt pares; in casibus reliquis (si vel utraque tres coëfficientes primos simul pares habet, vel neutra) aequivalebunt. --- Hoc modo adhuc multo plures propositiones speciales exhibere possemus, si supra (art.~277) plura exempla evoluta fuissent.

II. Pro omnibus hisce casibus poterit etiam, designantibus $f$, $f'$ formas ternarias aequivalentes, transformatio una alterius in alteram inveniri. Nam pro omnibus casibus in quavis classe formarum ternariarum multitudo satis parva formarum supra assignata est, ad quarum aliquam per methodos uniformes quaevis forma eiusdem classis reduci possit; has omnes ad unicam reducere ibidem docuimus. Sit $F$ haec forma in ea classe, in qua sunt $f$, $f'$, poteruntque per praecepta supra tradita inveniri transformationes formarum $f$, $f'$ in $F$, nec non formae $F$ in $f$, $f'$. Hinc per art.~270 deduci poterunt transformationes formae $f$ in $f'$ formaeque $f'$ in $f$.

III. Superesset itaque tantummodo, ostendere, quo pacto ex una transformatione formae ternariae $f$ in aliam $f'$ omnes transformationes possibiles derivari possint. Hoc problema pendet ab alio simpliciori, scilicet invenire omnes transformationes formae ternariae $f$ in se ipsam. Nimirum si $f$ per plures substitutiones $(\tau)$, $(\tau')$, $(\tau'')$ etc. in se ipsam et per substitutionem $(t)$ in $f'$ transit, patet si ad normam art.~270 combinetur transformatio $(t)$ cum $(\tau)$, $(\tau')$, $(\tau'')$ etc., prodire transformationes, per quas omnes $f$ in $f'$ transeat; praeterea per calculum facile probatur, quamvis transformationem formae $f$ in $f'$ hoc modo deduci posse e combinatione transformationis datae $(t)$ formae $f$ in $f'$ cum aliqua (et quidem
\SourceLeaf{338}
\textit{unica}) transformatione formae $f$ in se ipsam, adeoque ex combinatione transformationis datae formae $f$ in $f'$ cum \textit{omnibus} transformationibus formae $f$ in se ipsam oriri \textit{omnes} transformationes formae $f$ in $f'$, et quidem singulas semel tantum.

Investigationem omnium transformationum formae $f$ in se ipsam ad eum casum hic restringimus, ubi $f$ est forma definita, cuius coëfficientes $4,5,6$ omnes $=0$\HistoricalNote{*)}{Casus reliqui ubi $f$ est forma definita, ad hunc reduci possunt; si vero $f$ est forma indefinita, methodus omnino diversa adhibenda, transformationumque multitudo infinita erit.}. Sit itaque $f=\left(\begin{smallmatrix}a,\,a',\,a''\\0,\,0,\,0\end{smallmatrix}\right)$, exhibeanturque omnes substitutiones, per quas $f$ in se ipsam transit, indefinite per
\[
\begin{array}{ccc}
\alpha,&\beta,&\gamma\\
\alpha',&\beta',&\gamma'\\
\alpha'',&\beta'',&\gamma''
\end{array}
\]
ita ut satisfieri debeat aequationibus
\[
\begin{aligned}
a\alpha\alpha+a'\alpha'\alpha'+a''\alpha''\alpha''&=a&&\ldots\quad(\Omega)\\
a\beta\beta+a'\beta'\beta'+a''\beta''\beta''&=a'\\
a\gamma\gamma+a'\gamma'\gamma'+a''\gamma''\gamma''&=a''\\
a\alpha\beta+a'\alpha'\beta'+a''\alpha''\beta''&=0\\
a\alpha\gamma+a'\alpha'\gamma'+a''\alpha''\gamma''&=0\\
a\beta\gamma+a'\beta'\gamma'+a''\beta''\gamma''&=0
\end{aligned}
\]
Iam tres casus sunt distinguendi:

I. Quando $a$, $a'$, $a''$ (qui idem signum habebunt) omnes sunt inaequales, supponamus $a<a'$, $a'<a''$ (si alius magnitudinis ordo adest, eaedem conclusiones prorsus simili modo eruentur). Tunc aequ. prima in $(\Omega)$ manifesto requirit, ut sit $\alpha'=\alpha''=0$. adeoque $\alpha=\pm1$; hinc per aequ.~$4,5$ erit $\beta=0$. $\gamma=0$; similiter ex aequ.~$2$ erit $\beta''=0$, et proin $\beta'=\pm1$; hinc fit, per aequ.~$6$, $\gamma'=0$, et per $3$, $\gamma''=\pm1$, ita ut (ob signorum ambiguitatem independentem) omnino habeantur $8$ transformationes diversae.

II. Quando e numeris $a$, $a'$, $a''$ duo sunt aequales, \textit{e.~g.} $a'=a''$, tertius inaequalis, supponamus

\textit{primo} $a<a'$. Tunc eodem modo ut in casu praec. erit $\alpha'=0$, $\alpha''=0$, $\alpha=\pm1$, $\beta=0$, $\gamma=0$; ex aequ.~$2,3,6$ autem facile deducitur, esse debere
\SourceLeaf{339}
vel $\beta'=\pm1$, $\gamma'=0$, $\beta''=0$, $\gamma''=\pm1$, vel $\beta'=0$, $\gamma'=\pm1$, $\beta''=\pm1$, $\gamma''=0$.

Si vero, \textit{secundo}, $a>a'$, eaedem conclusiones sic obtinentur: ex aequ.~$2,3$ necessario erit $\beta=0$, $\gamma=0$, et vel $\beta'=\pm1$, $\gamma'=0$, $\beta''=0$, $\gamma''=\pm1$, vel $\beta'=0$, $\gamma'=\pm1$, $\beta''=\pm1$, $\gamma''=0$; pro suppositione utraque ex aequ.~$4,5$ erit $\alpha'=0$, $\alpha''=0$, atque ex $1$, $\alpha=\pm1$. Habentur itaque, pro utroque casu, $16$ transformationes diversae. --- Duo casus reliqui, ubi vel $a=a''$, vel $a=a'$, prorsus simili modo absolvuntur, si modo characteres $\alpha$, $\alpha'$, $\alpha''$ in priori cum $\beta$, $\beta'$, $\beta''$, in posteriori cum $\gamma$, $\gamma'$, $\gamma''$ resp. commutantur.

III. Quando omnes $a$, $a'$, $a''$ aequales sunt, aequationes $1,2,3$ requirunt, ut e tribus numeris $\alpha$, $\alpha'$, $\alpha''$, nec non ex $\beta$, $\beta'$, $\beta''$, ut et ex $\gamma$, $\gamma'$, $\gamma''$ bini sint $=0$, tertius $=\pm1$. Per aequ.~$4,5,6$ autem facile intelligitur, e tribus numeris $\alpha$, $\beta$, $\gamma$ unum tantummodo $=\pm1$ esse posse, similiterque ex $\alpha'$, $\beta'$, $\gamma'$, nec non ex $\alpha''$, $\beta''$, $\gamma''$. Quamobrem sex tantummodo combinationes dantur
\[
\begin{array}{c|c|c|c|c|c|c|}
\alpha&\alpha&\alpha'&\alpha'&\alpha''&\alpha''&=\pm1\\
\beta'&\beta''&\beta&\beta''&\beta&\beta'&=\pm1\\
\gamma''&\gamma'&\gamma''&\gamma&\gamma'&\gamma&=\pm1
\end{array}
\quad\text{Coëfficientes seni reliqui }=0
\]
ita ut ob signorum ambiguitatem omnino $48$ transformationes habeantur. --- Idem typus etiam casus praecedentes complectitur: sed e sex columnis primis prima sola accipi debet, quando $a$, $a'$, $a''$ omnes sunt inaequales; columna prima et secunda, quando $a'=a''$; prima et tertia, quando $a=a'$; prima et sexta, quando $a=a''$.

Hinc colligitur, si forma $f=axx+a'x'x'+a''x''x''$ in aliam aequivalentem $f'$ transeat per substitutionem
\[
x=\delta y+\varepsilon y'+\zeta y'',\quad x'=\delta'y+\varepsilon'y'+\zeta'y'',\quad x''=\delta''y+\varepsilon''y'+\zeta''y''
\]
omnes transf. formae $f$ in $f'$ contineri sub schemate sequente:
\[
\begin{array}{c|c|c|c|c|c|l}
x&x&x'&x'&x''&x''&=\pm(\delta y+\varepsilon y'+\zeta y'')\\
x'&x''&x&x''&x&x'&=\pm(\delta'y+\varepsilon'y'+\zeta'y'')\\
x''&x'&x''&x&x'&x&=\pm(\delta''y+\varepsilon''y'+\zeta''y'')
\end{array}
\]
eo discrimine, ut sex columnae primae omnes adhibendae sint, quando $a=a'=a''$;
\SourceLeaf{340}
columna $1$ et $2$, quando $a'$, $a''$ aequales, $a$ inaequalis; $1$ et $3$, quando $a=a'$; $1$ et $6$, quando $a=a''$; denique columna prima sola, quando $a$, $a'$, $a''$ omnes inaequales. In casu primo transformationum multitudo erit $48$, in secundo, tertio et quarto $16$, in quinto. $8$.

\Needspace{16\baselineskip}
\begin{center}
$\begin{array}{ccc}*&\quad&*\\&*&\end{array}$
\end{center}

\begin{center}
\textit{QUAEDAM APPLICATIONES AD THEORIAM FORMARUM BINARIARUM.}
\end{center}

\begin{center}
\textit{De invenienda forma, e cuius duplicatione forma binaria data generis principalis oriatur.}
\end{center}

Ab hac succincta primorum elementorum theoriae formarum ternariarum expositione ad quasdem applicationes speciales progredimur, inter quas primum locum meretur sequens

\Needspace{8\baselineskip}
\Article{286}
\textsc{Problema.} \textit{Proposita forma binaria $F=(A,B,C)$ determinantis $D$ ad genus principale pertinente: invenire formam binariam $f$, e cuius duplicatione illa oriatur.}

\textit{Sol.} I. Quaeratur repraesentatio propria formae ipsi $F$ oppositae $F'=ATT-2BTU+CUU$ per formam. ternariam $xx-2yz$, quae sit
\[
x=\alpha T+\beta U,\quad y=\alpha'T+\beta'U,\quad z=\alpha''T+\beta''U
\]
quod fieri posse e theoria praec. formarum ternariarum facile colligitur. Quum enim $F$ per hyp. sit e genere principali, dabitur valor expr. $\surd(A,\allowbreak B,\allowbreak C)(\mathrm{mod.}\,\allowbreak D)$, unde inveniri poterit forma ternaria $\varphi$ determinantis $1$, in quam $(A,-B,C)$ tamquam pars ingrediatur, cuius formae coëfficientes omnes fore integros nullo negotio perspicietur. Aeque facile intelligitur, $\varphi$ fore formam indefinitam (quoniam per hyp. $F$ certo non est forma negativa); unde necessario formae $xx-2yz$ aequivalens erit. Assignari poterit itaque transformatio huius in illam, quae repraesentationem propriam formae $F'$ per $xx-2yz$ suppeditabit. --- Tunc igitur erit
\[
A=\alpha\alpha-2\alpha'\alpha'',\quad -B=\alpha\beta-\alpha'\beta''-\alpha''\beta',\quad C=\beta\beta-2\beta'\beta''
\]
\SourceLeaf{341}
porro designatis numeris $\alpha\beta'-\alpha'\beta$, $\alpha'\beta''-\alpha''\beta'$, $\alpha''\beta-\alpha\beta''$ per $a,b,c$ resp., hi divisorem communem non habebunt, eritque $D=bb-2ac$.

II. Hinc adiumento observationis ultimae art.~235 facile concluditur, $F$ transire per substitutionem $2\beta',\allowbreak \beta,\allowbreak \beta,\allowbreak \beta'';\ 2\alpha',\allowbreak \alpha,\allowbreak \alpha,\allowbreak \alpha''$ in productum formae $(2a,-b,c)$ in se ipsam, nec non per substitutionem $\beta',\allowbreak \beta,\allowbreak \beta,\allowbreak 2\beta'';\ \alpha',\allowbreak \alpha,\allowbreak \alpha,\allowbreak 2\alpha''$ in productum formae $(a,-b,2c)$ in se ipsam. Iam divisor communis maximus numerorum $2a,2b,2c$ est $2$; si itaque $c$ est impar, $2a,2b,c$ divisorem communem non habebunt, sive $(2a,-b,c)$ erit forma proprie primitiva; similiter, si $a$ est impar, $(a,-b,2c)$ forma proprie primitiva erit; in casu priori $F$ oritur ex duplicatione formae $(2a,-b,c)$, in posteriori ex duplicatione formae $(a,-b,2c)$, (V. concl.~4, art.~235); unus vero horum casuum certo semper locum habebit. Si enim uterque $a,c$ esset par, $b$ necessario foret impar; iam facile confirmatur, esse $\beta''a+\beta b+\beta'c=0$, $\alpha''a+\alpha b+\alpha'c=0$, unde sequeretur, $\beta b$, $\alpha b$, adeoque etiam $\alpha$ et $\beta$ esse pares. Hinc autem $A$ et $C$ forent pares, quod esset contra hypothesin, secundum quam $F$ est forma e genere principali adeoque ex ordine proprie primitivo. --- Ceterum fieri etiam potest, ut tum $a$ tum $c$ impares sint, in quo itaque casu duae statim formae habebuntur, e quarum duplicatione $F$ oritur.

\textit{Ex.} Proposita sit forma $F=(5.\,2,31)$, det.~$-151$. Valor expressionis $\surd(5,2,31)$ hic invenitur $(55,22)$; hinc forma ternaria $\varphi=\left(\begin{smallmatrix}5,&31,&4\\11,&0,&-2\end{smallmatrix}\right)$; huic per praecepta art.~272 aequivalens invenitur forma $\left(\begin{smallmatrix}1,&1.&-1\\0,&0,&0\end{smallmatrix}\right)$, quae in $\varphi$ transit per substitutionem $\left\{\begin{smallmatrix}2,&2,&1\\1,&-6,&-2\\0,&3,&1\end{smallmatrix}\right\}$. Hinc adiumento transformationum in art.~277 traditarum invenitur, $\left(\begin{smallmatrix}1,&0,&0\\-1,&0,&0\end{smallmatrix}\right)$ transire in $\varphi$ per substitutionem $\left\{\begin{smallmatrix}3.&-7,&-2\\2,&-1,&0\\1,&-9,&-3\end{smallmatrix}\right\}$. Fit itaque $a=11$, $b=-17$, $c=20$; quare quum $a$ sit impar, $F$ oritur ex duplicatione formae $(11,17,40)$ transitque in productum huius formae in se ipsam per substitutionem $-1,-7,-7,-18;\ 2,3,3,2$.

\begin{center}
\textit{Omnibus characteribus, praeter eos, qui in artt.\ $263,264$ impossibiles inventi sunt, genera revera respondent.}
\end{center}

\Needspace{8\baselineskip}
\Article{287}
Circa problema in art. praec. solutum sequentes adhuc annotationes adiicimus.

I. Si forma $F$ per substitutionem $p,p',p'',p''';\ q,q',q'',q'''$ in productum e duabus formis $(h,i,k)$, $(h',i',k')$ transformatur, (utraque uti semper suppo%
